% ------------------------------------------------------------------ 3. C&C
\seccionnueva
\section{Estructura de componentes y conectores (C\&C)}
\label{sec:cc}

La vista C\&C muestra los elementos que existen \emph{en ejecución} y cómo se comunican. Las
flechas indican quién inicia la llamada o el acceso; HTTP y WebSocket transportan respuestas en
ambos sentidos. La figura~\ref{fig:cc} (página siguiente, en horizontal) es la vista completa y
las tablas posteriores son su catálogo.

\begin{table}[H]
\centering
\small
\filas
\begin{tabularx}{\linewidth}{@{}P{3.4cm} P{1.1cm} L@{}}
  \encabezado \hd{Tipo de elemento} & \hd{Cant.} & \hd{Elementos} \\
  Presentación & 2 & P1 Web SSR (Next.js, TypeScript) · P2 Cliente local (Electron, TypeScript) \\
  Integración & 2 & I1 API Gateway (Nginx) · I2 RabbitMQ (broker AMQP) \\
  Lógica de negocio & 7 & L1 Auth (Go) · L2 Colaboración (Python) · L3 Pizarra (Python) · L4 Chat (Java) · L5 Pixels (Python) · L6 NotificationAPI (C\#) · L7 MessageWorker (C\#) \\
  Datos & 5 & D1 UserDB · D2 BoardDB · D3 CanvasDB · D5 NotificationDB (PostgreSQL) · D4 PixelDB con la base de Chat (MongoDB) \\
  Conectores & 12 & C01–C12: HTTP, WebSocket, HTTP interno, AMQP, PostgreSQL y MongoDB \\
  \bottomrule
\end{tabularx}
\caption{Resumen de la vista C\&C.}
\label{tab:resumen-cc}
\end{table}

\subsection*{Cómo leer la vista}

\begin{itemize}
  \item \textbf{De arriba hacia abajo}: presentación, gateway, lógica de negocio, asincronía y datos.
  \item \textbf{Ningún cliente toca una base de datos} ni una API directamente: todo pasa por el gateway.
  \item \textbf{Cada almacén tiene un único dueño}; los demás servicios le preguntan por HTTP interno (C05) o le envían eventos (C06).
  \item \textbf{La notificación es asíncrona}: si RabbitMQ o el \textit{worker} están caídos, la operación de negocio igual se confirma y el evento espera en el outbox.
  \item \textbf{La difusión WebSocket de Pizarra y Chat es local a cada proceso}: una réplica por servicio de salas; RabbitMQ no la convierte en un sistema de salas distribuido.
\end{itemize}

% Vista en página horizontal para que el diagrama sea legible.
\begin{landscape}
\begin{figure}[p]
\centering
\resizebox{\linewidth}{!}{%
\begin{tikzpicture}[x=1cm,y=1cm]
  \tikzset{api/.style={comp=#1,minimum width=3.7cm,text width=3.4cm}}

  % --- Presentación
  \node[client,minimum width=4.6cm] (browser) at (-8.4,0.4)
    {\textbf{Navegador}\\[-1pt]{\scriptsize HTML del servidor + JS hidratado}};
  \node[comp=cbViolet,minimum width=4.6cm,text width=4.3cm] (web) at (-8.4,-2.5)
    {\textbf{P1 · Web SSR}\\[-1pt]Next.js · TypeScript\\[-1pt]{\scriptsize SSR por petición, BFF, Server Actions}};
  \node[comp=cbViolet,minimum width=4.8cm,text width=4.5cm] (desktop) at (8.4,0.4)
    {\textbf{P2 · Cliente local}\\[-1pt]Electron · TypeScript\\[-1pt]{\scriptsize renderer empaquetado, IPC acotado}};

  % --- Integración
  \node[comp=cbAmber,minimum width=8.2cm,text width=7.9cm] (gw) at (0,-2.5)
    {\textbf{I1 · API Gateway} · Nginx\\[-1pt]{\scriptsize único acceso público · rutas REST y proxy WebSocket · X-Request-Id · bloquea /internal}};

  % --- Lógica de negocio
  \node[api=cbIndigo] (auth)   at (-10.5,-5.4) {\textbf{L1 · Auth}\\[-1pt]Go\\[-1pt]{\scriptsize registro, login, JWT}};
  \node[api=cbIndigo] (collab) at (-6.3,-5.4)  {\textbf{L2 · Colaboración}\\[-1pt]Python\\[-1pt]{\scriptsize tableros, roles, invitaciones}};
  \node[api=cbIndigo] (canvas) at (-2.1,-5.4)  {\textbf{L3 · Pizarra}\\[-1pt]Python\\[-1pt]{\scriptsize salas, ops y versiones}};
  \node[api=cbIndigo] (chat)   at (2.1,-5.4)   {\textbf{L4 · Chat}\\[-1pt]Java\\[-1pt]{\scriptsize historial y mensajes}};
  \node[api=cbIndigo] (pixels) at (6.3,-5.4)   {\textbf{L5 · Pixels}\\[-1pt]Python\\[-1pt]{\scriptsize lienzo 200×200, cooldown}};
  \node[api=cbIndigo] (notif)  at (10.5,-5.4)  {\textbf{L6 · NotificationAPI}\\[-1pt]C\#\\[-1pt]{\scriptsize eventos, outbox, bandeja}};

  % --- Asincronía
  \node[broker,minimum width=3.2cm] (mq) at (6.9,-9.1) {\textbf{I2 · RabbitMQ}\\[-1pt]{\scriptsize cola durable · DLQ}};
  \node[api=cbIndigo] (worker) at (10.8,-9.1) {\textbf{L7 · MessageWorker}\\[-1pt]C\#\\[-1pt]{\scriptsize consumo idempotente}};

  % --- Datos
  \node[db] (userdb)  at (-10.5,-11.7) {\textbf{D1 · UserDB}\\PostgreSQL};
  \node[db] (boarddb) at (-6.3,-11.7)  {\textbf{D2 · BoardDB}\\PostgreSQL};
  \node[db] (canvasdb) at (-2.1,-11.7) {\textbf{D3 · CanvasDB}\\PostgreSQL+JSONB};
  \node[db,minimum width=26mm] (mongo) at (2.9,-11.7) {\textbf{D4 · PixelDB}\\MongoDB\\{\tiny + base \texttt{chat}}};
  \node[db] (notifdb) at (10.8,-11.7)  {\textbf{D5 · NotificationDB}\\PostgreSQL};

  % --- C01..C03: presentación
  \draw[http] (browser) -- node[lbl,left] {C01 HTTP: documento, BFF} (web);
  \draw[http] (web) -- node[lbl,above] {C02 HTTP interno} (gw);
  \draw[ws] (browser.south east) -- node[lbl,sloped,above] {C03 WebSocket (cookie)} (gw.north west);
  \draw[http] (desktop.south west) -- node[lbl,sloped,above] {C03 HTTP + WS (Bearer)} (gw.north east);

  % --- C04: gateway -> APIs (bus)
  \coordinate (bus) at (0,-3.95);
  \draw[thick,cbInk!75] (gw.south) -- (bus);
  \draw[thick,cbInk!75] (-10.5,-3.95) -- (10.5,-3.95);
  \foreach \n/\x in {auth/-10.5,collab/-6.3,canvas/-2.1,chat/2.1,pixels/6.3,notif/10.5}
    \draw[http] (\x,-3.95) -- (\n.north);
  \fill[cbInk!75] (bus) circle (1.6pt);
  \node[lbl,anchor=south] at (-4.2,-3.95) {C04 HTTP REST a las seis APIs · WS a Pizarra, Chat y Pixels};

  % --- C05: permisos (Pizarra y Chat -> Colaboración)
  \draw[internal] (1.2,-6.12) -- (1.2,-6.75) -- (-5.4,-6.75) -- (-5.4,-6.12);
  \draw[internal,-] (-1.2,-6.12) -- (-1.2,-6.75);
  \fill[cbCyan!80!black] (-1.2,-6.75) circle (1.6pt);
  \node[lbl,text=cbCyan!60!black] at (-3.3,-6.75) {C05 HTTP: autorización};

  % --- C06: eventos (Colaboración y Chat -> NotificationAPI)
  \draw[internal] (-7.2,-6.12) -- (-7.2,-7.45) -- (9.9,-7.45) -- (9.9,-6.12);
  \draw[internal,-] (3.0,-6.12) -- (3.0,-7.45);
  \fill[cbCyan!80!black] (3.0,-7.45) circle (1.6pt);
  \node[lbl,text=cbCyan!60!black] at (-1.6,-7.45) {C06 HTTP interno: eventos pendientes con reintento};

  % --- C07: asincronía
  \draw[amqp] (notif.south west) -- node[lbl,sloped,above] {C07 AMQP: publicar} (mq.north);
  \draw[amqp] (mq) -- node[lbl,below=2pt] {consumir} (worker);

  % --- C08..C12: datos
  \draw[sql] (auth) -- node[lbl] {C08 SQL} (userdb);
  \draw[sql] (collab) -- node[lbl] {C09 SQL} (boarddb);
  \draw[sql] (canvas) -- node[lbl] {C10 SQL} (canvasdb);
  \draw[sql] (chat.south) -- node[lbl,pos=0.45] {C11} (mongo.north west);
  \draw[sql] (pixels.south) -- node[lbl,pos=0.38] {C11 Mongo} (mongo.north east);
  \draw[sql] (worker) -- node[lbl] {C12 SQL} (notifdb);
  \draw[sql] (notif.east) -- ++(1.55,0) |- node[lbl,pos=0.25] {C12} (notifdb.east);

  % --- Leyenda
  \node[draw=cbLine,fill=white,rounded corners=3pt,anchor=north,inner sep=6pt,font=\sffamily\scriptsize]
    at (0,1.15) {%
    \begin{tabular}{@{}l@{\hspace{4pt}}l@{\hspace{14pt}}l@{\hspace{4pt}}l@{}}
      \tikz\draw[http] (0,0)--(0.8,0); & HTTP / REST & \tikz\draw[amqp] (0,0)--(0.8,0); & AMQP (RabbitMQ) \\
      \tikz\draw[ws] (0,0)--(0.8,0); & WebSocket & \tikz\draw[sql] (0,0)--(0.8,0); & Conector de datos \\
      \tikz\draw[internal] (0,0)--(0.8,0); & HTTP interno & & \\
    \end{tabular}};
\end{tikzpicture}}
\caption{Vista C\&C del prototipo 2: dos presentaciones, gateway, seis APIs de negocio y un
\textit{worker}, broker de mensajes y cinco almacenes de datos (dos motores, cinco lenguajes).}
\label{fig:cc}
\end{figure}
\end{landscape}

\subsection{Catálogo de componentes}

\begin{longtable}{@{}P{2.9cm} P{3.2cm} P{6.3cm} P{2.3cm}@{}}
  \encabezado \hd{Componente} & \hd{Tecnología} & \hd{Responsabilidad} & \hd{Estado} \\
  \endfirsthead
  \encabezado \hd{Componente} & \hd{Tecnología} & \hd{Responsabilidad} & \hd{Estado} \\
  \endhead
  P1 · Web SSR & Next.js 16, React 19, TypeScript & Renderiza por petición el dashboard y la cabecera de cada tablero; hidrata lienzo, chat y bandeja; BFF hacia el gateway & \hecho \\
  \rowcolor{cbGray} P2 · Cliente local & Electron, React, TypeScript & Presentación de escritorio que reutiliza \code{core/} y \code{ui/} de la web con su propio adaptador de transporte & \porhacer \\
  I1 · API Gateway & Nginx & Único punto de entrada: enruta REST, reenvía WebSocket, propaga \code{X-Request-Id}, no expone rutas internas & \integracion \\
  \rowcolor{cbGray} L1 · Auth & Go (chi), Argon2id, JWT HS256 & Registro, login y \code{/auth/me}; emite el JWT que usan todos los servicios & \pone \\
  L2 · Colaboración & Python (FastAPI) & Único dueño de tableros, miembros, roles e invitaciones; endpoint interno de autorización & \integracion \\
  \rowcolor{cbGray} L3 · Pizarra & Python (FastAPI, WebSocket) & Salas de dibujo con orden \code{seq}, \code{ack}/\code{saved}, presencia, autosave e historial & \integracion \\
  L4 · Chat & Java & Historial paginado por cursor, envío idempotente y difusión a miembros autorizados & \porhacer \\
  \rowcolor{cbGray} L5 · Pixels & Python (FastAPI) & Lienzo colectivo 200×200 con enfriamiento por usuario & \pone \\
  L6 · NotificationAPI & C\# (.NET) & Recibe eventos internos, guarda evento y outbox en una transacción, publica en RabbitMQ y expone la bandeja & \porhacer \\
  \rowcolor{cbGray} L7 · MessageWorker & C\# (.NET) & Consume tareas, crea notificaciones de forma idempotente y confirma tras persistir & \porhacer \\
  I2 · RabbitMQ & AMQP 0-9-1 & Cola durable con confirmaciones de publicación, \textit{ack} manual y cola de mensajes fallidos & \porhacer \\
  \bottomrule
  \caption{Componentes lógicos y de integración.}
  \label{tab:componentes}
\end{longtable}

\begin{longtable}{@{}P{3.3cm} P{2.7cm} P{3.6cm} P{5.1cm}@{}}
  \encabezado \hd{Almacén} & \hd{Motor} & \hd{Único usuario} & \hd{Contenido} \\
  \endfirsthead
  \encabezado \hd{Almacén} & \hd{Motor} & \hd{Único usuario} & \hd{Contenido} \\
  \endhead
  D1 · UserDB & PostgreSQL 16 & Auth & Usuarios y hashes Argon2id \\
  \rowcolor{cbGray} D2 · BoardDB & PostgreSQL 16 & Colaboración & \code{boards}, \code{board\_members}, \code{board\_invites} \\
  D3 · CanvasDB & PostgreSQL 16 + JSONB & Pizarra & \code{board\_states}, \code{board\_versions}; \code{board\_id} como referencia externa (sin FK) \\
  \rowcolor{cbGray} D4 · PixelDB & MongoDB 7 & Pixels; Chat en su base \code{chat} & Píxeles, enfriamientos y actividad; mensajes aislados por base y credencial \\
  D5 · NotificationDB & PostgreSQL 16 & NotificationAPI y MessageWorker & Eventos, outbox, notificaciones con unicidad \code{(event\_id, recipient\_id)} \\
  \bottomrule
  \caption{Componentes de datos: cinco instancias, dos motores, sin claves foráneas entre bases.}
  \label{tab:datos}
\end{longtable}

\subsection{Catálogo de conectores}

\begin{longtable}{@{}P{1.1cm} P{4.4cm} P{2.9cm} P{6.3cm}@{}}
  \encabezado \hd{ID} & \hd{Origen → destino} & \hd{Protocolo} & \hd{Propósito} \\
  \endfirsthead
  \encabezado \hd{ID} & \hd{Origen → destino} & \hd{Protocolo} & \hd{Propósito} \\
  \endhead
  C01 & Navegador → Web SSR & HTTP & Documento SSR, recursos, BFF (\code{/bff/*}) y Server Actions; publicado por el gateway en el mismo origen \\
  \rowcolor{cbGray} C02 & Web SSR → Gateway & HTTP interno & Consultas autorizadas durante el SSR y reenvío del BFF con \code{Authorization: Bearer} \\
  C03 & Navegador / cliente local → Gateway & HTTP REST y WebSocket & Acciones y tiempo real; el navegador autentica el WS con la cookie, Electron con cabecera \\
  \rowcolor{cbGray} C04 & Gateway → seis APIs & HTTP REST; WS a Pizarra, Chat y Pixels & Enrutamiento por prefijo de ruta pública \\
  C05 & Pizarra / Chat → Colaboración & HTTP interno & Resolver el rol del usuario sobre el tablero antes de aplicar operaciones o entregar mensajes \\
  \rowcolor{cbGray} C06 & Colaboración / Chat → NotificationAPI & HTTP interno & Eventos \code{invite.accepted} y \code{chat.message.created} con \code{event\_id} estable y reintentos \\
  C07 & NotificationAPI → RabbitMQ → Worker & AMQP & Publicación confirmada y consumo con \textit{ack} manual \\
  \rowcolor{cbGray} C08 & Auth → UserDB & PostgreSQL & Cuentas \\
  C09 & Colaboración → BoardDB & PostgreSQL & Tableros y membresías \\
  \rowcolor{cbGray} C10 & Pizarra → CanvasDB & PostgreSQL & Estado y versiones \\
  C11 & Pixels / Chat → MongoDB & Protocolo MongoDB & Bases y credenciales separadas \\
  \rowcolor{cbGray} C12 & NotificationAPI / Worker → NotificationDB & PostgreSQL & Outbox, notificaciones y lecturas \\
  \bottomrule
  \caption{Conectores de la vista C\&C.}
  \label{tab:conectores}
\end{longtable}

\subsection{Estilos y patrones arquitectónicos}

\begin{table}[H]
\centering
\small
\filas
\begin{tabularx}{\linewidth}{@{}P{3.9cm} L@{}}
  \encabezado \hd{Estilo o patrón} & \hd{Dónde y por qué} \\
  Cliente-servidor & Ambas presentaciones son clientes de servicios remotos; ningún cliente accede a una base de datos \\
  Servicios distribuidos & Cada dominio (identidad, colaboración, dibujo, chat, píxeles, notificaciones) es un proceso y contenedor independiente con su propio almacén \\
  API Gateway & Un solo ingreso público: simplifica clientes, centraliza enrutamiento y correlación, y oculta las rutas internas \\
  Backend for Frontend & La web SSR traduce la cookie de sesión en \code{Bearer} y aplica la protección CSRF en un solo lugar \\
  SSR + hidratación & El HTML inicial trae los datos autorizados; el navegador solo agrega interacción \\
  Publicador-suscriptor (WS) & Pizarra, Chat y Pixels difunden eventos a todos los conectados de una sala \\
  Productor-consumidor con cola & NotificationAPI produce, MessageWorker consume; el trabajo sobrevive a reinicios del \textit{worker} \\
  Outbox transaccional & El evento se guarda con el cambio de negocio y se publica después; una caída del broker no pierde notificaciones \\
  Consumidor idempotente & Unicidad \code{(event\_id, recipient\_id)} y \code{client\_message\_id}: reintentos sin duplicados visibles \\
  Persistencia políglota & PostgreSQL para datos relacionales y JSONB; MongoDB para escrituras frecuentes de píxeles y mensajes \\
  \bottomrule
\end{tabularx}
\caption{Estilos y patrones presentes en la vista C\&C.}
\label{tab:estilos}
\end{table}

\begin{decision}{ajustes al boceto inicial del C\&C}
\begin{itemize}
  \item \textbf{CanvasDB sigue en PostgreSQL con JSONB}: el código ya tiene autosave, versiones e historial; migrarlo a MongoDB no lo exige el enunciado y añadiría riesgo.
  \item \textbf{PixelDB sigue en MongoDB} y Chat usa una base lógica propia en esa instancia: no se cuenta como otro contenedor.
  \item \textbf{NotificationAPI también lee NotificationDB}: la bandeja y la marca de lectura la necesitan; el \textit{worker} escribe las notificaciones.
  \item \textbf{Pizarra y Chat consultan permisos a Colaboración por HTTP} y no tocan BoardDB.
  \item \textbf{La web usa un BFF}: las APIs reciben siempre \code{Bearer}, igual desde la web que desde Electron.
\end{itemize}
\end{decision}

\begin{pendiente}{Marlon (vista C\&C y conectores) · Cristian y David (eventos)}
Confirmar contra el código final: rutas reales del gateway, protocolo del WebSocket de Chat,
nombres de colas y DLQ, y quién emite cada evento. Actualizar la columna \emph{Estado}
de la tabla~\ref{tab:componentes}.
\end{pendiente}
